% Part: first-order-logic
% Chapter: natural-deduction
% Section: provability-propositional

% Vérification des propriétés de dérivabilité utilisées pour les ensembles
% maximalement cohérents dans le chapitre sur la complétude.

\documentclass[../../../include/open-logic-section]{subfiles}

\begin{document}

\iftag{FOL}
      {\olfileid[fr]{fol}{ntd}{ppr}}
      {\olfileid[fr]{pl}{ntd}{ppr}}

\olsection{\usetoken{S}{derivability} et connecteurs propositionnels}

\begin{explain}
Nous montrons que la relation de !!{derivability}~$\Proves$ en
déduction naturelle permet d'établir quelques faits élémentaires
concernant les connecteurs propositionnels, notamment
$!A \land !B \Proves !A$ et $!A, !A \lif !B \Proves !B$
(modus ponens). Ces faits seront utilisés dans la démonstration
du théorème de complétude.
\end{explain}

\begin{prop}\ollabel{prop:provability-land}
  \begin{enumerate}
  \item \ollabel{prop:provability-land-left} On a à la fois
    $!A \land !B \Proves !A$ et $!A \land !B \Proves !B$.
  \item \ollabel{prop:provability-land-right} $!A, !B \Proves !A \land !B$.
  \end{enumerate}
\end{prop}

\begin{proof}
  \begin{enumerate}
  \item Voici les deux !!{derivation}s :
\begin{prooftree}
      \AxiomC{$!A \land !B$}
      \RightLabel{\Elim{\land}}
      \UnaryInfC{$!A$}
      \DisplayProof\qquad\bottomAlignProof
      \AxiomC{$!A \land !B$}
      \RightLabel{\Elim{\land}}
      \UnaryInfC{$!B$}
    \end{prooftree}
\item Nous avons la !!{derivation} suivante :
\begin{prooftree}
      \AxiomC{$!A$}
      \AxiomC{$!B$}
      \RightLabel{\Intro{\land}}
      \BinaryInfC{$!A \land !B$}
    \end{prooftree}
\end{enumerate}
\end{proof}

\begin{prop}\ollabel{prop:provability-lor}
  \begin{enumerate}
  \item $!A \lor !B, \lnot !A, \lnot !B$ est incohérent.
  \item On a à la fois $!A \Proves !A \lor !B$ et
    $!B \Proves !A \lor !B$.
  \end{enumerate}
\end{prop}

\begin{proof}
  \begin{enumerate}
  \item Considérons la !!{derivation} suivante :
\begin{prooftree}
      \AxiomC{$!A \lor !B$}
      \AxiomC{$\lnot !A$}
      \AxiomC{$\Discharge{!A}{1}$}
      \RightLabel{\Elim{\lnot}}
      \BinaryInfC{$\lfalse$}
      \AxiomC{$\lnot !B$}
      \AxiomC{$\Discharge{!B}{1}$}
      \RightLabel{\Elim{\lnot}}
      \BinaryInfC{$\lfalse$}
      \DischargeRule{\Elim{\lor}}{1}
      \TrinaryInfC{$\lfalse$}
    \end{prooftree}
C'est une !!{derivation} de~$\lfalse$ dont les hypothèses
    !!{undischarged}s sont $!A \lor !B$, $\lnot !A$ et $\lnot !B$.
  \item Voici les deux !!{derivation}s :
\begin{prooftree}
      \AxiomC{$!A$}
      \RightLabel{\Intro{\lor}}
      \UnaryInfC{$!A \lor !B$}
      \DisplayProof\qquad\bottomAlignProof
      \AxiomC{$!B$}
      \RightLabel{\Intro{\lor}}
      \UnaryInfC{$!A \lor !B$}
    \end{prooftree}
\end{enumerate}
\end{proof}

\begin{prop}\ollabel{prop:provability-lif}
  \begin{enumerate}
  \item \ollabel{prop:provability-lif-left} $!A, !A \lif !B \Proves !B$.
  \item \ollabel{prop:provability-lif-right}
    On a à la fois $\lnot !A \Proves !A \lif !B$ et
    $!B \Proves !A \lif !B$.
  \end{enumerate}
\end{prop}

\begin{proof}
  \begin{enumerate}
  \item Nous avons la !!{derivation} suivante :
\begin{prooftree}
      \AxiomC{$!A \lif !B$}
      \AxiomC{$!A$}
      \RightLabel{\Elim{\lif}}
      \BinaryInfC{$!B$}
    \end{prooftree}
\item Les deux !!{derivation}s suivantes établissent le résultat :
\begin{prooftree}
      \AxiomC{$\lnot !A$}
      \AxiomC{$\Discharge{!A}{1}$}
      \RightLabel{\Elim{\lnot}}
      \BinaryInfC{$\lfalse$}
      \RightLabel{\FalseInt}
      \UnaryInfC{$!B$}
      \DischargeRule{\Intro{\lif}}{1}
      \UnaryInfC{$!A \lif !B$}
      \DisplayProof\qquad\bottomAlignProof
      \AxiomC{$!B$}
      \RightLabel{\Intro{\lif}}
      \UnaryInfC{$!A \lif !B$}
    \end{prooftree}
Remarquons que $\Intro{\lif}$ permet de !!{discharge}
    l'hypothèse~$!A$, sans l'exiger.
  \end{enumerate}
\end{proof}

\end{document}
